\chapter{Plan del Proyecto}
\label{anx:plan}

\section{Justificación y viabilidad}

\subsection{Justificación}

La investigación es conveniente porque ataca directamente el cuello de botella
descrito en el Capítulo~\ref{cap:generalidades}: el análisis manual del audio
recolectado. El monitoreo acústico pasivo con sensores de bajo costo permite
recolectar a gran escala, pero ese potencial queda desaprovechado mientras el
análisis siga siendo lento, costoso y laborioso. Un sistema de pre-anotación
asistida hace que el monitoreo a gran escala resulte operativamente viable y
permite que el tiempo experto se dedique a la interpretación ecológica en lugar
de al procesamiento.

En el plano práctico, el proyecto reduce la infrautilización del material acústico
acumulado y, con ella, la latencia entre la recolección y la generación de
conocimiento. En el plano teórico, aborda dos vacíos confirmados de forma
independiente por la revisión del Capítulo~\ref{cap:estado}: la ausencia de
modelos y de conjuntos de datos con anotación tiempo--frecuencia para primates
neotropicales, y la falta de estrategias establecidas para el etiquetado
jerárquico de especie y tipo de llamada sobre eventos detectados.

\subsection{Beneficios esperados}

\begin{itemize}[nosep]
    \item Reducir la carga de revisión por hora de audio analizada: menos
          material que inspeccionar para la misma cobertura de eventos.
    \item Poner en uso el audio ya grabado que hoy permanece sin anotar.
    \item Acortar la latencia entre la recolección y la obtención de resultados
          ecológicos.
    \item Disminuir el costo y el error humano asociados al análisis manual.
    \item Entregar una herramienta reutilizable, con código y pesos públicos,
          integrable al flujo de trabajo existente en Raven.
\end{itemize}

\subsection{Viabilidad técnica}

El proyecto no depende de ningún desarrollo previo inexistente. Las
arquitecturas necesarias están disponibles y documentadas
\citep{he_deep_2016,kahl_birdnet_2021,nolasco_learning_2023}; las herramientas
son software libre de uso extendido (Tabla~\ref{tab:herramientas}); los
recursos de cómputo se cubren con acceso a GPU en la nube; los datos iniciales,
aproximadamente 10\,GB de grabaciones WAV capturadas con AudioMoth, con sus
anotaciones de Raven, ya están entregados; y la metodología de trabajo
(CRISP-DM) está establecida y es la que estructura el documento.

\subsection{Viabilidad temporal}

El cronograma es factible dentro del calendario académico, con dos puntos de
tensión identificados. El primero es OE1: la curación y estandarización del
conjunto es meticulosa y cualquier retraso ahí se propaga a todo lo que sigue.
El segundo es OE2, que consume tiempo tanto de cómputo como de análisis humano
en la experimentación. Juega a favor que no haya recolección de datos en campo:
el material de partida ya existe.

\subsection{Viabilidad económica}

El grueso del software es de código abierto y sin costo. Los sensores ya están
desplegados y su costo no recae sobre este proyecto. El desembolso real se
concentra en cómputo en la nube y en servicios de almacenamiento; el resto del
presupuesto es una valorización del tiempo dedicado, no un pago efectivo.

\section{Definición y alcance del proyecto}

\subsection{Alcance}

El proyecto desarrolla y evalúa un sistema de pre-anotación asistida que produce
cajas tiempo--frecuencia etiquetadas sobre grabaciones de monitoreo acústico
pasivo, para las ocho especies de la Tabla~\ref{tab:especies}. El sistema debe
cumplir tres tareas: localizar el evento en tiempo y frecuencia, identificar la
especie y clasificar el tipo de vocalización. Una parte central del alcance es
la caracterización cuantitativa de su comportamiento operativo: el tiempo de
revisión que cuesta una hora de audio con y sin la asistencia del sistema, la
carga de propuestas que genera y la naturaleza de sus discrepancias. La medición
cronometrada y la revisión experta de una muestra de predicciones están dentro
del alcance y dependen de la disponibilidad del equipo de investigación, que es
el riesgo R3 de la sección~\ref{sec:riesgos}. El proyecto
recorre el ciclo completo de CRISP-DM, desde la preparación de los datos hasta el
procesamiento de grabaciones nuevas.

\subsection{Entregables y criterios de aceptación}

Los entregables son los resultados esperados enumerados en la
sección~\ref{sec:resultados}. La Tabla~\ref{tab:aceptacion} recoge el criterio
de aceptación de cada uno: la condición verificable que permite declararlo
terminado.

\begin{table}[htbp]
    \centering
    \caption{Criterios de aceptación}
    \label{tab:aceptacion}
    \small
    \begin{tabular}{@{}L{1.5cm}L{12.9cm}@{}}
        \toprule
        \textbf{Entreg.} & \textbf{Criterio de aceptación}                                                                                                                                                                                  \\
        \midrule
        RE1.1            & El protocolo define reglas de sintaxis para las etiquetas de especie y tipo de llamada, y especifica el tratamiento de errores, duplicados e inconsistencias de las anotaciones originales.                      \\
        \addlinespace
        RE1.2            & Todas las anotaciones del conjunto final cumplen las reglas de RE1.1, y se entrega un manifiesto con la asignación de cada archivo de grabación a una partición.                                                 \\
        \addlinespace
        RE1.3            & La ficha describe la estructura de directorios, incluye un diccionario de datos por columna y declara transformaciones, sesgos y limitaciones conocidas.                                                         \\
        \addlinespace
        RE2.1            & El documento justifica la arquitectura propuesta y la línea base a partir del estado del arte, y fija la configuración e hiperparámetros iniciales de ambas.                                                     \\
        \addlinespace
        RE2.2            & El \textit{pipeline} ejecuta el ciclo completo de entrenamiento, evaluación e inferencia, e incluye instrucciones de instalación y ejecución.                                                                    \\
        \addlinespace
        RE2.3            & El informe presenta las métricas de los tres ejes para la propuesta y la línea base, discute los resultados y concluye con una recomendación justificada.                                                        \\
        \addlinespace
        RE2.4            & Los pesos del modelo se cargan sin error y un script de demostración produce predicciones en el formato de tabla de selección de Raven sobre un clip de prueba.                                                  \\
        \addlinespace
        RE2.5            & El repositorio público contiene el código del \textit{pipeline}, los pesos del modelo final y las instrucciones para reproducir la inferencia.                                                                   \\
        \addlinespace
        RE3.1            & La metodología, fechada antes de la evaluación, especifica el subconjunto de prueba, el punto de operación (confianza 0,5 y NMS IoU 0,3), el umbral IoU 0,5 que define un verdadero positivo, las métricas, el emparejamiento y el procedimiento de revisión de discrepancias. \\
        \addlinespace
        RE3.2            & Se reportan los minutos por hora de audio en condición manual y asistida sobre material comparable, con el orden de revisión declarado, y el tiempo de procesamiento del modelo por hora de audio.                                                  \\
        \addlinespace
        RE3.3            & El informe clasifica la muestra de predicciones de alta confianza en evento no anotado, ruido y error de encuadre, e incluye ejemplos sobre el espectrograma de aciertos, de errores comunes y de detecciones correctas ausentes de la referencia.  \\
        \bottomrule
    \end{tabular}
\end{table}

\subsection{Exclusiones}

\begin{itemize}[nosep]
    \item No se desarrolla un sistema de monitoreo en tiempo real desplegable en
          campo; el alcance es el procesamiento de grabaciones existentes.
    \item El análisis detallado de sonidos antropogénicos queda fuera del foco
          primario, más allá de su detección y categorización general.
    \item El estudio se limita a las ocho especies declaradas.
    \item No se desarrolla hardware de adquisición.
    \item No se realiza nueva recolección de datos en campo.
    \item \textbf{No se cubre la distribución ni el mantenimiento del conjunto de
          datos ni del modelo.} Las grabaciones y las anotaciones primarias
          pertenecen al equipo de investigación, y publicarlas, licenciarlas o
          sostenerlas como recurso de terceros son decisiones suyas y no de esta
          tesis. El repositorio que compromete RE2.5 contiene el código y los
          artefactos del modelo final; empaquetar, versionar para terceros o dar
          soporte queda fuera.
\end{itemize}

\subsection{Limitaciones y restricciones}

\rotulo{Limitaciones.} La ausencia de conjuntos etiquetados a gran escala para
primates neotropicales obliga a invertir un esfuerzo considerable en OE1, y la
calidad de las anotaciones manuales recibidas es subjetiva y variable. El
entorno acústico amazónico es denso y presenta con frecuencia baja relación
señal-ruido. La variabilidad intraespecífica de las vocalizaciones y su
naturaleza graduada dificultan definir categorías discretas, y las llamadas
superpuestas o de baja intensidad son especialmente difíciles de modelar. Los
sensores de bajo costo añaden variabilidad entre dispositivos. Por último, sin
revisión experta las predicciones de alta confianza situadas en regiones sin
anotación (tipo D de la Tabla~\ref{tab:taxonomia-fp}) no pueden clasificarse de
forma concluyente entre eventos omitidos por la referencia y falsos positivos
genuinos: esta tesis las caracteriza geométricamente y las reporta como una
cantidad acotada, y deja su resolución como trabajo futuro condicionado a la
disponibilidad del equipo de investigación.

\rotulo{Restricciones.} El desarrollo está sujeto al calendario académico. El
proyecto depende del conjunto de datos entregado y de la verificación de su
calidad. El acceso a GPU limita la escala de la experimentación, en particular
la búsqueda exhaustiva de hiperparámetros. El desarrollo, la implementación y la
evaluación recaen sobre el autor, con la supervisión del asesor.

\section{Planificación del proyecto}

\subsection{Estructura de desglose del trabajo}

La estructura de desglose se organiza según las fases de CRISP-DM y los
objetivos específicos. La Tabla~\ref{tab:edt} resume los paquetes de trabajo y
el entregable asociado a cada uno, y la Figura~\ref{fig:wbs-tree} los dibuja en
jerarquía: sólo se abren los tres paquetes que producen resultados esperados.

\begin{table}[htbp]
    \centering
    \caption{EDT del proyecto}
    \label{tab:edt}
    \small
    \begin{tabular}{@{}L{0.8cm}L{4.4cm}L{6.6cm}L{2.2cm}@{}}
        \toprule
        \textbf{ID} & \textbf{Paquete}                 & \textbf{Contenido}                                                                                 & \textbf{Entregable} \\
        \midrule
        1           & Gestión del proyecto             & Planificación, seguimiento, comunicación con el asesor, gestión de riesgos y de la calidad, cierre & N/A                 \\
        \addlinespace
        2           & Comprensión de la investigación  & Objetivos, formulación del problema de detección y criterios de éxito                              & N/A                 \\
        \addlinespace
        3           & Comprensión de los datos         & Acceso al corpus, descripción del formato, análisis exploratorio y verificación de calidad         & N/A                 \\
        \addlinespace
        4.1         & Limpieza y estandarización       & Corrección de etiquetas, tratamiento de audios dañados                                             & RE1.1               \\
        4.2         & Curación del conjunto            & Aplicación del protocolo, vocabulario cerrado, partición por grabación                             & RE1.2               \\
        4.3         & Ingeniería de características    & Preprocesamiento, representación tiempo--frecuencia, ventaneo, aumento de datos                    & N/A                 \\
        4.4         & Documentación del conjunto       & Ficha del conjunto derivado                                                                        & RE1.3               \\
        \addlinespace
        5.1         & Diseño del detector              & Arquitectura propuesta y línea base, con justificación                                             & RE2.1               \\
        5.2         & Implementación                   & \textit{Pipeline} de entrenamiento, evaluación e inferencia                                        & RE2.2               \\
        5.3         & Publicación del código           & Repositorio con código y pesos                                                                     & RE2.5               \\
        \addlinespace
        6.1         & Evaluación cuantitativa          & Métricas por eje, comparación con la línea base, análisis de errores                               & RE2.3               \\
        6.2         & Selección del modelo final       & Modelo elegido y exportador a tabla de selección de Raven                                          & RE2.4               \\
        6.3         & Metodología de validación        & Subconjunto de prueba, punto de operación, emparejamiento, métricas y protocolo de revisión experta & RE3.1               \\
        6.4         & Medición del tiempo de revisión  & Cronometraje de la revisión manual y asistida sobre material comparable, y costo de inferencia      & RE3.2               \\
        6.5         & Análisis de errores              & Clasificación de las predicciones de alta confianza y ejemplos sobre el espectrograma               & RE3.3               \\
        6.6         & Revisión experta de la muestra   & Resolución con el equipo de investigación de las predicciones sin anotación                         & RE3.2, RE3.3        \\
        \addlinespace
        7           & Despliegue                       & Inferencia por lotes sobre grabaciones no vistas y generación de tablas de anotación               & N/A                 \\
        \addlinespace
        8           & Documentación y difusión         & Redacción de la tesis y preparación de la defensa                                                  & N/A                 \\
        \bottomrule
    \end{tabular}
\end{table}

\begin{figure}[htbp]
    \centering
    \begin{tikzpicture}[
            x=1mm, y=1mm,
            paquete/.style={caja, draw=azul!70, fill=azul!6, align=left, font=\footnotesize,
                    text width=58mm, minimum height=5.2mm, inner sep=2.5pt, anchor=north west},
            abierto/.style={paquete, draw=azul!80!black, fill=azul, text=white},
            hoja/.style={caja, draw=black!40, fill=white, align=left, font=\scriptsize,
                    text width=59mm, minimum height=5.2mm, inner sep=2.5pt, anchor=north west},
            rama/.style={draw=black!45},
        ]
        \node[paquete] (p1) at (0,    0)      {\textbf{1.} Gestión del proyecto};
        \node[paquete] (p2) at (0,   -6.4)    {\textbf{2.} Comprensión de la investigación};
        \node[paquete] (p3) at (0,  -12.8)    {\textbf{3.} Comprensión de los datos};
        \node[abierto] (p4) at (0,  -19.2)    {\textbf{4.} Preparación de datos};
        \node[abierto] (p5) at (0,  -51.2)    {\textbf{5.} Modelado};
        \node[abierto] (p6) at (0,  -76.8)    {\textbf{6.} Evaluación};
        \node[paquete] (p7) at (0, -121.6)    {\textbf{7.} Despliegue};
        \node[paquete] (p8) at (0, -128.0)    {\textbf{8.} Documentación y difusión};
        \foreach \y/\t/\r in {%
                -25.6/{\textbf{4.1} Limpieza y estandarización}/{RE1.1},
                -32.0/{\textbf{4.2} Curación del conjunto}/{RE1.2},
                -38.4/{\textbf{4.3} Ingeniería de características}/{},
                -44.8/{\textbf{4.4} Documentación del conjunto}/{RE1.3},
                -57.6/{\textbf{5.1} Diseño del detector}/{RE2.1},
                -64.0/{\textbf{5.2} Implementación}/{RE2.2},
                -70.4/{\textbf{5.3} Publicación del código}/{RE2.5},
                -83.2/{\textbf{6.1} Evaluación cuantitativa}/{RE2.3},
                -89.6/{\textbf{6.2} Selección del modelo final}/{RE2.4},
                -96.0/{\textbf{6.3} Metodología de validación}/{RE3.1},
                -102.4/{\textbf{6.4} Medición del tiempo de revisión}/{RE3.2},
                -108.8/{\textbf{6.5} Análisis de errores}/{RE3.3},
                -115.2/{\textbf{6.6} Revisión experta}/{RE3.2, RE3.3}} {
                \node[hoja] at (64, \y)
                {\t\hfill\textcolor{verde!55!black}{\bfseries\r}};
                \draw[rama] (4, \y-2.6) -- (64, \y-2.6);
            }
        \draw[rama] (4, -24.4) -- (4, -47.4);
        \draw[rama] (4, -56.4) -- (4, -73.0);
        \draw[rama] (4, -82.0) -- (4, -117.8);
    \end{tikzpicture}
    \caption{EDT en árbol}
    \label{fig:wbs-tree}
\end{figure}

\subsection{Cronograma}

El cronograma se ajusta al calendario académico y se organiza en bloques que
respetan las dependencias entre paquetes: la preparación de datos precede al
modelado, el modelado precede a la evaluación, y la fijación del punto de operación
precede a la medición del tiempo y al análisis de errores, porque ambos se
realizan sobre las predicciones del modelo ya seleccionado en ese punto. La
sesión de revisión con el equipo de investigación se agenda con holgura, porque
alimenta a la vez la condición asistida de RE3.2 y la muestra resuelta de RE3.3.


\subsection{Estimación de costos}

Los costos se presentan en dólares estadounidenses y distinguen desembolso real
de valorización. Las Tablas~\ref{tab:costos-hw} a~\ref{tab:costos-ind} detallan
cada categoría, y la Tabla~\ref{tab:costos-resumen} las consolida.

\begin{table}[htbp]
    \centering
    \small
    \begin{threeparttable}
        \caption{Costos de equipos}
        \label{tab:costos-hw}
        \begin{tabular}{@{}L{4.2cm}L{6.4cm}r@{}}
            \toprule
            \textbf{Categoría}     & \textbf{Descripción}                                               & \thead{USD}        \\
            \midrule
            Almacenamiento externo & SSD de 512\,GB para el conjunto de datos y sus derivados           & 150,00             \\
            \addlinespace
            Equipo de desarrollo   & Laptop con GPU dedicada\tnote{a}                                   & 1\,500,00          \\
            \addlinespace
            Cómputo en la nube     & Instancia con GPU de 48\,GiB de memoria, estimada por horas de uso & 600,00             \\
            \midrule
            \textbf{Subtotal}      &                                                                    & \textbf{2\,250,00} \\
            \bottomrule
        \end{tabular}
        \begin{tablenotes}[flushleft]
            \footnotesize
            \item[a] Equipo preexistente propiedad del tesista. No representa un
            desembolso del proyecto, pero sí un recurso valorizado.
        \end{tablenotes}
    \end{threeparttable}
\end{table}

\begin{table}[htbp]
    \centering
    \small
    \begin{threeparttable}
        \caption{Capital intelectual}
        \label{tab:costos-rrhh}
        \begin{tabular}{@{}L{5.0cm}rrr@{}}
            \toprule
            \textbf{Recurso}         & \thead{Horas} & \thead{USD/h} & \thead{Total USD}  \\
            \midrule
            Tesista\tnote{a}         & 600           & 5,00          & 2\,500,00          \\
            Asesor de tesis\tnote{b} & 40            & 25,00         & 1\,000,00          \\
            \midrule
            \textbf{Subtotal}        &               &               & \textbf{3\,500,00} \\
            \bottomrule
        \end{tabular}
        \begin{tablenotes}[flushleft]
            \footnotesize
            \item[a] Valorización del tiempo dedicado; no implica un pago real.
            \item[b] Cubierto por la institución.
        \end{tablenotes}
    \end{threeparttable}
\end{table}

\begin{table}[htbp]
    \centering
    \caption{Costos de software y servicios}
    \label{tab:costos-sw}
    \small
    \begin{tabular}{@{}L{5.0cm}L{3.2cm}rrr@{}}
        \toprule
        \textbf{Software o servicio}      & \textbf{Licencia}      & \thead{Meses} & \thead{USD/mes} & \thead{Total USD} \\
        \midrule
        Python, uv, Pandas, NumPy         & Código abierto         & 12            & 0,00            & 0,00              \\
        PyTorch, torchaudio, Transformers & Código abierto         & 12            & 0,00            & 0,00              \\
        Visual Studio Code                & Código abierto         & 12            & 0,00            & 0,00              \\
        Zotero                            & Gratuito               & 12            & 0,00            & 0,00              \\
        Raven Lite                        & Licencia de estudiante & 12            & 0,00            & 0,00              \\
        Entorno de cómputo con GPU        & Suscripción            & 12            & 9,99            & 119,88            \\
        Almacenamiento en la nube         & Suscripción            & 12            & 1,67            & 20,00             \\
        \midrule
        \textbf{Subtotal}                 &                        &               &                 & \textbf{139,88}   \\
        \bottomrule
    \end{tabular}
\end{table}

\begin{table}[htbp]
    \centering
    \caption{Costos indirectos}
    \label{tab:costos-ind}
    \small
    \begin{tabular}{@{}L{5.4cm}L{4.6cm}rr@{}}
        \toprule
        \textbf{Categoría}  & \textbf{Descripción}                           & \thead{Meses} & \thead{Total USD} \\
        \midrule
        Internet            & Conexión para investigación y acceso a la nube & 12            & 480,00            \\
        Electricidad        & Consumo del equipo y periféricos               & 12            & 360,00            \\
        Impresiones y otros & Materiales de estudio e impresiones            & 12            & 120,00            \\
        \midrule
        \textbf{Subtotal}   &                                                &               & \textbf{960,00}   \\
        \bottomrule
    \end{tabular}
\end{table}

\begin{table}[htbp]
    \centering
    \caption{Resumen de costos}
    \label{tab:costos-resumen}
    \small
    \begin{tabular}{@{}L{7.0cm}r@{}}
        \toprule
        \textbf{Categoría}              & \thead{Total USD}  \\
        \midrule
        Equipamiento (hardware)         & 2\,250,00          \\
        Recursos humanos (valorización) & 3\,500,00          \\
        Software y servicios            & 139,88             \\
        Costos indirectos               & 960,00             \\
        \midrule
        \textbf{Total}                  & \textbf{6\,849,88} \\
        \bottomrule
    \end{tabular}
\end{table}

\subsection{Lista de recursos}

\rotulo{Personas.} El tesista, responsable del desarrollo, la implementación y
el análisis; y el asesor, que aporta guía, supervisión y validación experta.

\rotulo{Equipamiento.} Equipo de desarrollo con GPU dedicada, almacenamiento
externo para el conjunto de datos y sus derivados, y acceso a servidores con GPU
para el entrenamiento.

\rotulo{Herramientas.} Las de la Tabla~\ref{tab:herramientas}, más el acceso a
Raven para la validación cualitativa de las predicciones.

\subsection{Gestión de riesgos}
\label{sec:riesgos}

La Tabla~\ref{tab:riesgos} recoge los riesgos identificados con su probabilidad,
impacto, severidad y las acciones previstas.

\begin{table}[htbp]
    \centering
    \caption{Riesgos del proyecto}
    \label{tab:riesgos}
    \small
    \begin{tabular}{@{}L{2.8cm}ccL{4.4cm}L{4.2cm}@{}}
        \toprule
        \textbf{Riesgo}                                                       & \thead{Prob.} & \thead{Sever.} & \textbf{Mitigación} & \textbf{Contingencia} \\
        \midrule
        R1. Calidad de los datos: anotaciones inconsistentes y audios dañados &
        Baja                                                                  & Media         &
        Validación cruzada de una muestra con el asesor; detección automática de
        archivos corruptos o vacíos                                           &
        Excluir el material bajo umbral mínimo de calidad y documentar la
        inconsistencia como limitación                                                                                                                       \\
        \addlinespace
        R2. Recursos de cómputo insuficientes                                 &
        Baja                                                                  & Media         &
        Espectrogramas precalculados; \textit{pipeline} de datos optimizado;
        planificación de las corridas largas                                  &
        Reducir la complejidad de las configuraciones menos prometedoras;
        acotar la búsqueda de hiperparámetros; parada temprana                                                                                               \\
        \addlinespace
        R3. Baja disponibilidad de revisores expertos para la validación      &
        Media                                                                 & Baja          &
        Sesión de revisión agendada con holgura y con el material ya preparado; la
        muestra se dimensiona para que una sola sesión la resuelva               &
        Reducir el tamaño de la muestra revisada y declarar de forma explícita
        la limitación de alcance de RE3.2 y RE3.3                                                                                                            \\
        \addlinespace
        R4. Falla del equipo de desarrollo                                    &
        Alta                                                                  & Media         &
        Código, datos y documentos sincronizados en repositorio remoto; entorno
        replicable por archivo de configuración                               &
        Reinstalar el sistema, clonar el repositorio y recrear el entorno desde
        su especificación                                                                                                                                    \\
        \addlinespace
        R5. Incompatibilidad del entorno de ejecución                         &
        Media                                                                 & Media         &
        Requerimientos explícitos y versionados; contenedor de ejecución;
        versiones estables de las bibliotecas clave                           &
        Proveer un entorno preconfigurado listo para ejecutar; migrar a
        versiones compatibles                                                                                                                                \\
        \addlinespace
        R6. Aumento del costo de cómputo en la nube                           &
        Baja                                                                  & Baja          &
        Desarrollo y depuración en recursos locales o gratuitos; uso eficiente
        de la GPU; priorización de experimentos                               &
        Reducir el alcance de la experimentación; solicitar créditos académicos;
        usar recursos de la universidad                                                                                                                      \\
        \addlinespace
        R7. Retrasos por disponibilidad limitada del tesista                  &
        Alta                                                                  & Alta          &
        Cronograma con holgura basado en la EDT; bloques de tiempo protegidos;
        seguimiento semanal del avance                                        &
        Replanificar al detectar la desviación; negociar reducción de alcance
        con el asesor; priorizar la ruta crítica                                                                                                             \\
        \bottomrule
    \end{tabular}
\end{table}


% ===================================================================
% REFERENCIAS BIBLIOGRÁFICAS
% ===================================================================
